% ================================================================
% MINI-ÉVALUATION — ECNSA
% Structure calquée sur minievaluation.tex du Chap2.
%
% UTILISATION :
%   1. Copier ce fichier dans le dossier du devoir.
%   2. Remplacer TOUS les champs [entre crochets].
%   3. Compiler avec pdflatex (pas de bibliographie = une seule passe suffit).
%
% LIMITE : ~5 pages. Pas de fancyhdr, numéro de page en bas centré.
% ================================================================
\documentclass[11pt,a4paper]{article}

% Chargement du préambule — cherche le fichier dans plusieurs chemins
% relatifs pour fonctionner depuis n'importe quel sous-dossier du projet.
\newcommand{\chargerpreambule}[1]{%
  \IfFileExists{../../templates/#1}{\input{../../templates/#1}}{%
  \IfFileExists{../templates/#1}{\input{../templates/#1}}{%
  \IfFileExists{templates/#1}{\input{templates/#1}}{\input{#1}}}}}
\chargerpreambule{preamble-commun.tex}
\chargerpreambule{preamble-mini-evaluation.tex}

% ----------------------------------------------------------------
% MÉTADONNÉES — à personnaliser
% ----------------------------------------------------------------
\title{\textbf{Mini-évaluation pratique\\{}[Titre du devoir]}}
\author{[Prénom NOM]\\ Classe : [classe]}
\date{[date]}     % ou \date{\today} pour la date du jour

\begin{document}

  \maketitle
  \thispagestyle{empty}   % pas de numéro sur la page de titre

  % ================================================================
  % SECTION 1 — Première question ou premier thème
  % Renommer le titre selon l'énoncé.
  % ================================================================
  \section{[Besoin / Question 1]}

  [Contexte ou énoncé de la question.]

  \subsection{[Sous-partie 1]}

  \begin{itemize}
    \item \textbf{[Point clé :]} [explication].
    \item \textbf{[Point clé :]} [explication].
  \end{itemize}

  \subsection{[Sous-partie 2]}

  \begin{itemize}
    \item \textbf{[Point clé :]} [explication].
  \end{itemize}

  % Dupliquer les blocs \subsection pour chaque sous-partie supplémentaire.

  % \newpage   % décommenter si la section suivante doit commencer sur une nouvelle page

  % ================================================================
  % SECTION 2 — Deuxième question ou deuxième thème
  % ================================================================
  \section{[Règles de politique / Question 2]}

  \subsection{[Règle 1] --- [Intitulé]}

  \begin{itemize}
    \item \textbf{Énoncé :} [Décrire la règle ou l'exigence.]
    \item \textbf{Périmètre :} [Systèmes, réseaux ou actifs concernés.]
    \item \textbf{Mesure technique :} [Action à mettre en œuvre.]
    \item \textbf{Preuve de contrôle :} [Comment vérifier que la mesure est en place.]
  \end{itemize}

  \subsection{[Règle 2] --- [Intitulé]}

  \begin{itemize}
    \item \textbf{Énoncé :} [Décrire la règle ou l'exigence.]
    \item \textbf{Périmètre :} [Systèmes, réseaux ou actifs concernés.]
    \item \textbf{Mesure technique :} [Action à mettre en œuvre.]
    \item \textbf{Preuve de contrôle :} [Comment vérifier que la mesure est en place.]
  \end{itemize}

  % Dupliquer le bloc \subsection ci-dessus pour chaque règle supplémentaire.

\end{document}
